\kbvsechead{(8)\hspace{0.45em}ಅವಧಾರಣಾರ್ಥಕಾವ್ಯಯ :}

\indent ನಿಶ್ಚಯಾರ್ಥದಲ್ಲಿ ಬರುವ ಅವ್ಯಯಗಳಿಗೆ ಅವಧಾರಣಾರ್ಥಕಾವ್ಯಯಗಳೆಂದು ಹೆಸರು.\par

\vspace{3pt plus 2pt}\noindent\hspace*{1.2em}\kbvlabel{ಉದಾ}\par\vspace{1pt}

\begin{kbvderiv}
\vspace{4pt plus 3pt minus 1pt}
\begin{center}\begin{adjustbox}{max width=\textwidth}
{\renewcommand{\arraystretch}{1.22}\setlength{\tabcolsep}{4pt}
\begin{tabular}{l c l}
ಅದೇ ಶಿವಗಂಗೆ (ಅದು + ಏ & {\color{kbvgrey}$=$} & \kbvresult{ಅದೇ)} \\
ನಾನೇ ಕಣ್ಣಾರೆ ನೋಡಿದೆ. ...(ನಾನು + ಏ & {\color{kbvgrey}$=$} & \kbvresult{ನಾನೇ)} \\
ಶಿವನೇ ನನ್ನ ರಕ್ಷಕ ...(ಶಿವನು + ಏ & {\color{kbvgrey}$=$} & \kbvresult{ಶಿವನೇ)} \\
ಅವಳೇ ಭಾಗ್ಯಶಾಲಿ (ಅವಳು + ಏ & {\color{kbvgrey}$=$} & \kbvresult{ಅವಳೇ)} \\
\end{tabular}}
\end{adjustbox}\end{center}
\vspace{4pt plus 3pt minus 1pt}
\end{kbvderiv}

\indent ಇಲ್ಲಿ ಅದೇ, ನಾನೇ, ಶಿವನೇ, ಅವಳೇ ಎಂಬ ಪದಗಳಲ್ಲಿ "ಏ" ಕಾರಣವಾಗಿ, ಅವಧಾರಣಾರ್ಥಕಾವ್ಯಯ. ಈ ಪದಗಳ ಪ್ರಯೋಗದಲ್ಲಿ ನಿಶ್ಚಯ ಆಗುವುದು. ನಿರ್ಣಯ ಕಂಡುಬರುವುದು.\par

\kbvexercise{ಅಭ್ಯಾಸ ಪತ್ರಿಕೆ ೧೬}

\vspace{2.5pt plus 3pt minus 0.5pt}\noindent\hspace*{1.2em}\kbvnum{1.}\hspace{0.55em}ಅವ್ಯಯ ಎಂದರೇನು ? ಉದಾಹರಣೆಗಳ ಸಹಿತ ತಿಳಿಸಿ.\par

\vspace{2.5pt plus 3pt minus 0.5pt}\noindent\hspace*{1.2em}\kbvnum{2.}\hspace{0.55em}“ಹುಡುಗರು ಸುಮ್ಮನೆ ಬಂದರು” — ಈ ವಾಕ್ಯದಲ್ಲಿ ವಚನಗಳನ್ನು ಬದಲಾಯಿಸಿ ಬರೆಯಿರಿ.\par

\vspace{2.5pt plus 3pt minus 0.5pt}\noindent\hspace*{1.2em}\kbvnum{3.}\hspace{0.55em}ಅವ್ಯಯದಲ್ಲಿ ಎಷ್ಟು ಪ್ರಕಾರಗಳಿವೆ ? ಅವು ಯಾವುವು ?\par

\vspace{2.5pt plus 3pt minus 0.5pt}\noindent\hspace*{1.2em}\kbvnum{4.}\hspace{0.55em}ಇಲ್ಲಿ ಹಲವು ವಾಕ್ಯಗಳನ್ನು ಕೊಟ್ಟಿದೆ. ಪ್ರತಿಯೊಂದ ರಲ್ಲೂ ಯಾವ ಪ್ರಕಾರದ ಅವ್ಯಯವಿರುವುದನ್ನು ತಿಳಿದು ಹೇಳಿ.\par

\vspace{2.5pt plus 3pt minus 0.5pt}\noindent\hspace*{1.2em}\kbvnum{(1)}\hspace{0.55em}ದ್ರಾಕ್ಷಿಗೊಂಚಲು ಸುಲಭವಾಗಿ ಕೈಗೆಟುಕುತ್ತದೆ.\par

\vspace{2.5pt plus 3pt minus 0.5pt}\noindent\hspace*{1.2em}\kbvnum{(2)}\hspace{0.55em}ಗಾಡಿಯು ತಟ್ಟನೆ ನಿಂತಿತು.\par

\vspace{2.5pt plus 3pt minus 0.5pt}\noindent\hspace*{1.2em}\kbvnum{(3)}\hspace{0.55em}ಗಾಳಿಯು ಭರ್ರನೆ ಬೀಸಿತು.\par

\vspace{2.5pt plus 3pt minus 0.5pt}\noindent\hspace*{1.2em}\kbvnum{(4)}\hspace{0.55em}ಅಯ್ಯೋ! ಮಗು ರೈಲಿಗೆ ಸಿಕ್ಕಿತು.\par

\vspace{2.5pt plus 3pt minus 0.5pt}\noindent\hspace*{1.2em}\kbvnum{(5)}\hspace{0.55em}ಅವನಿಂದ ಹಣ ಬರಬೇಕು.\par

\vspace{2.5pt plus 3pt minus 0.5pt}\noindent\hspace*{1.2em}\kbvnum{(6)}\hspace{0.55em}ರಾಮ ಮತ್ತು ಕೃಷ್ಣ ಅವತಾರ ಪುರುಷರು.\par

\vspace{2.5pt plus 3pt minus 0.5pt}\noindent\hspace*{1.2em}\kbvnum{(7)}\hspace{0.55em}ನನಗೆ ಬೇಕಾಗಿರುವುದು ಕಾಫಿ ಅಥವಾ ಟೀ.\par

\vspace{2.5pt plus 3pt minus 0.5pt}\noindent\hspace*{1.2em}\kbvnum{(8)}\hspace{0.55em}ನೂಲಿನಂತೆ ಸೀರೆ ತಾಯಿಯಂತೆ ಮಗಳು.\par
